\section{Introduction}

Point process models are a standard tool for studying the temporal
structure of earthquake occurrences and quantifying aftershock activity.
The simplest benchmark is a homogeneous Poisson process, which assumes
independent exponentially distributed inter-event times. However,
observed catalogues typically exhibit strong temporal clustering, with
bursts of activity following large events, which motivates the use of
self-exciting Hawkes processes.

In this short note we compare five models on an earthquake catalogue
from the USGS Earthquake Hazards Program: a homogeneous Poisson process,
unmarked and marked exponential and power-law Hawkes processes. Our goal
is to identify a parsimonious specification that captures the observed
clustering and magnitude dependence of aftershock activity, as assessed
by residual diagnostics under the time-rescaling theorem.
